%!TEX program = xelatex
\documentclass[11pt,a4paper]{article}
\usepackage[margin=2.2cm]{geometry}
\usepackage{amsmath,amssymb,amsthm,mathtools}
\usepackage{bm}
\usepackage{enumitem}
\usepackage{booktabs}
\usepackage{xcolor}
\usepackage{hyperref}
\usepackage{fancyhdr}
\usepackage{xeCJK}

\hypersetup{
    colorlinks=true,
    linkcolor=blue!60!black,
    citecolor=blue!60!black,
    urlcolor=blue!60!black
}

\pagestyle{fancy}
\fancyhf{}
\rhead{\small Qiuzhen QE Probability \& Statistics --- Fall 2023}
\lhead{\small Official Qualifying Exam}
\rfoot{\small Page \thepage}

\DeclareMathOperator*{\argmax}{arg\,max}
\DeclareMathOperator*{\argmin}{arg\,min}
\DeclareMathOperator{\Var}{Var}
\DeclareMathOperator{\E}{\mathbb{E}}
\DeclareMathOperator{\Pbb}{\mathbb{P}}
\DeclareMathOperator{\Cov}{Cov}

\setlist[itemize]{topsep=3pt,itemsep=2pt,parsep=0pt}
\setlist[enumerate]{topsep=3pt,itemsep=2pt,parsep=0pt}

\newcommand{\examstatement}[1]{%
  \par\addvspace{0.85em}%
  \noindent{\bfseries #1}\par\nobreak\addvspace{0.28em}%
}

% Shared amsthm note/remark/theorem styles
% Shared math extras for qe-review.
% Requires already-loaded: amsmath, amssymb.
%
% Classic pitfall: AMS blackboard-bold fonts have no digit glyphs, so
% \mathbb{1} renders as overlapping garbage. Map the digit 1 to dsfont's
% \mathds{1}; leave \mathbb{R}, \mathbb{P}, etc. on AMS.
%
% Usage (path relative to the main .tex being compiled):
%   notes:  % Shared math extras for qe-review.
% Requires already-loaded: amsmath, amssymb.
%
% Classic pitfall: AMS blackboard-bold fonts have no digit glyphs, so
% \mathbb{1} renders as overlapping garbage. Map the digit 1 to dsfont's
% \mathds{1}; leave \mathbb{R}, \mathbb{P}, etc. on AMS.
%
% Usage (path relative to the main .tex being compiled):
%   notes:  % Shared math extras for qe-review.
% Requires already-loaded: amsmath, amssymb.
%
% Classic pitfall: AMS blackboard-bold fonts have no digit glyphs, so
% \mathbb{1} renders as overlapping garbage. Map the digit 1 to dsfont's
% \mathds{1}; leave \mathbb{R}, \mathbb{P}, etc. on AMS.
%
% Usage (path relative to the main .tex being compiled):
%   notes:  \input{../../common/qe-math-extras.tex}
%   exams:  \input{../../../common/qe-math-extras.tex}

\usepackage{dsfont}
\usepackage{pdftexcmds}

\makeatletter
\let\qe@amsmmathbb\mathbb
\renewcommand{\mathbb}[1]{%
  \ifnum\pdf@strcmp{\unexpanded{#1}}{1}=\z@
    \mathds{1}%
  \else
    \qe@amsmmathbb{#1}%
  \fi
}
\makeatother

% Preferred indicator macros (either style is safe after the remap above).
\newcommand{\bbone}{\mathds{1}}
\newcommand{\ind}{\mathbf{1}}

%   exams:  % Shared math extras for qe-review.
% Requires already-loaded: amsmath, amssymb.
%
% Classic pitfall: AMS blackboard-bold fonts have no digit glyphs, so
% \mathbb{1} renders as overlapping garbage. Map the digit 1 to dsfont's
% \mathds{1}; leave \mathbb{R}, \mathbb{P}, etc. on AMS.
%
% Usage (path relative to the main .tex being compiled):
%   notes:  \input{../../common/qe-math-extras.tex}
%   exams:  \input{../../../common/qe-math-extras.tex}

\usepackage{dsfont}
\usepackage{pdftexcmds}

\makeatletter
\let\qe@amsmmathbb\mathbb
\renewcommand{\mathbb}[1]{%
  \ifnum\pdf@strcmp{\unexpanded{#1}}{1}=\z@
    \mathds{1}%
  \else
    \qe@amsmmathbb{#1}%
  \fi
}
\makeatother

% Preferred indicator macros (either style is safe after the remap above).
\newcommand{\bbone}{\mathds{1}}
\newcommand{\ind}{\mathbf{1}}


\usepackage{dsfont}
\usepackage{pdftexcmds}

\makeatletter
\let\qe@amsmmathbb\mathbb
\renewcommand{\mathbb}[1]{%
  \ifnum\pdf@strcmp{\unexpanded{#1}}{1}=\z@
    \mathds{1}%
  \else
    \qe@amsmmathbb{#1}%
  \fi
}
\makeatother

% Preferred indicator macros (either style is safe after the remap above).
\newcommand{\bbone}{\mathds{1}}
\newcommand{\ind}{\mathbf{1}}

%   exams:  % Shared math extras for qe-review.
% Requires already-loaded: amsmath, amssymb.
%
% Classic pitfall: AMS blackboard-bold fonts have no digit glyphs, so
% \mathbb{1} renders as overlapping garbage. Map the digit 1 to dsfont's
% \mathds{1}; leave \mathbb{R}, \mathbb{P}, etc. on AMS.
%
% Usage (path relative to the main .tex being compiled):
%   notes:  % Shared math extras for qe-review.
% Requires already-loaded: amsmath, amssymb.
%
% Classic pitfall: AMS blackboard-bold fonts have no digit glyphs, so
% \mathbb{1} renders as overlapping garbage. Map the digit 1 to dsfont's
% \mathds{1}; leave \mathbb{R}, \mathbb{P}, etc. on AMS.
%
% Usage (path relative to the main .tex being compiled):
%   notes:  \input{../../common/qe-math-extras.tex}
%   exams:  \input{../../../common/qe-math-extras.tex}

\usepackage{dsfont}
\usepackage{pdftexcmds}

\makeatletter
\let\qe@amsmmathbb\mathbb
\renewcommand{\mathbb}[1]{%
  \ifnum\pdf@strcmp{\unexpanded{#1}}{1}=\z@
    \mathds{1}%
  \else
    \qe@amsmmathbb{#1}%
  \fi
}
\makeatother

% Preferred indicator macros (either style is safe after the remap above).
\newcommand{\bbone}{\mathds{1}}
\newcommand{\ind}{\mathbf{1}}

%   exams:  % Shared math extras for qe-review.
% Requires already-loaded: amsmath, amssymb.
%
% Classic pitfall: AMS blackboard-bold fonts have no digit glyphs, so
% \mathbb{1} renders as overlapping garbage. Map the digit 1 to dsfont's
% \mathds{1}; leave \mathbb{R}, \mathbb{P}, etc. on AMS.
%
% Usage (path relative to the main .tex being compiled):
%   notes:  \input{../../common/qe-math-extras.tex}
%   exams:  \input{../../../common/qe-math-extras.tex}

\usepackage{dsfont}
\usepackage{pdftexcmds}

\makeatletter
\let\qe@amsmmathbb\mathbb
\renewcommand{\mathbb}[1]{%
  \ifnum\pdf@strcmp{\unexpanded{#1}}{1}=\z@
    \mathds{1}%
  \else
    \qe@amsmmathbb{#1}%
  \fi
}
\makeatother

% Preferred indicator macros (either style is safe after the remap above).
\newcommand{\bbone}{\mathds{1}}
\newcommand{\ind}{\mathbf{1}}


\usepackage{dsfont}
\usepackage{pdftexcmds}

\makeatletter
\let\qe@amsmmathbb\mathbb
\renewcommand{\mathbb}[1]{%
  \ifnum\pdf@strcmp{\unexpanded{#1}}{1}=\z@
    \mathds{1}%
  \else
    \qe@amsmmathbb{#1}%
  \fi
}
\makeatother

% Preferred indicator macros (either style is safe after the remap above).
\newcommand{\bbone}{\mathds{1}}
\newcommand{\ind}{\mathbf{1}}


\usepackage{dsfont}
\usepackage{pdftexcmds}

\makeatletter
\let\qe@amsmmathbb\mathbb
\renewcommand{\mathbb}[1]{%
  \ifnum\pdf@strcmp{\unexpanded{#1}}{1}=\z@
    \mathds{1}%
  \else
    \qe@amsmmathbb{#1}%
  \fi
}
\makeatother

% Preferred indicator macros (either style is safe after the remap above).
\newcommand{\bbone}{\mathds{1}}
\newcommand{\ind}{\mathbf{1}}

% Shared amsthm environments for qe-review (minimal).
% Requires already-loaded: amsthm.
% Safe to \input after \usepackage{amsmath,amssymb,amsthm,...}.
%
% Shared numbering uses aliascnt so cleveref keys off the environment
% name (Lemma, Proposition, Exercise, ...) rather than the parent
% counter (Theorem, Definition, Remark). Without aliascnt, \cref{lem:...}
% prints ``Theorem 9.13'' instead of ``Lemma 9.13''.

\usepackage{aliascnt}

\theoremstyle{plain}
\newtheorem{theorem}{Theorem}[section]

\newaliascnt{lemma}{theorem}
\newtheorem{lemma}[lemma]{Lemma}
\aliascntresetthe{lemma}

\newaliascnt{proposition}{theorem}
\newtheorem{proposition}[proposition]{Proposition}
\aliascntresetthe{proposition}

\newaliascnt{corollary}{theorem}
\newtheorem{corollary}[corollary]{Corollary}
\aliascntresetthe{corollary}

\newaliascnt{claim}{theorem}
\newtheorem{claim}[claim]{Claim}
\aliascntresetthe{claim}

\newaliascnt{fact}{theorem}
\newtheorem{fact}[fact]{Fact}
\aliascntresetthe{fact}

\theoremstyle{definition}
\newtheorem{definition}{Definition}[section]

\newaliascnt{exercise}{definition}
\newtheorem{exercise}[exercise]{Exercise}
\aliascntresetthe{exercise}

\newaliascnt{assumption}{definition}
\newtheorem{assumption}[assumption]{Assumption}
\aliascntresetthe{assumption}

\newaliascnt{notation}{definition}
\newtheorem{notation}[notation]{Notation}
\aliascntresetthe{notation}

\newaliascnt{construction}{definition}
\newtheorem{construction}[construction]{Construction}
\aliascntresetthe{construction}

% Example: document-wide numbering (not per section / chapter)
\newtheorem{example}{Example}

\theoremstyle{remark}
\newtheorem{remark}{Remark}[section]
\newtheorem*{remark*}{Remark}

\newaliascnt{heuristic}{remark}
\newtheorem{heuristic}[heuristic]{Heuristic}
\aliascntresetthe{heuristic}

\newaliascnt{convention}{remark}
\newtheorem{convention}[convention]{Convention}
\aliascntresetthe{convention}

\newaliascnt{warning}{remark}
\newtheorem{warning}[warning]{Warning}
\aliascntresetthe{warning}

\newaliascnt{proofsketch}{remark}
\newtheorem{proofsketch}[proofsketch]{Proof sketch}
\aliascntresetthe{proofsketch}

\newaliascnt{checkpoint}{remark}
\newtheorem{checkpoint}[checkpoint]{Checkpoint}
\aliascntresetthe{checkpoint}

% Note: document-wide numbering (not per section)
\newtheorem{note}{Note}
\newtheorem*{note*}{Note}

\newaliascnt{examnote}{note}
\newtheorem{examnote}[examnote]{Exam note}
\aliascntresetthe{examnote}

\newaliascnt{study}{note}
\newtheorem{study}[study]{Study note}
\aliascntresetthe{study}

\newaliascnt{summary}{note}
\newtheorem{summary}[summary]{Summary}
\aliascntresetthe{summary}

\newtheorem{examproblem}{Problem}[section]
\newtheorem*{examproblem*}{Problem}

% Bilingual aliases
\newenvironment{dingli}{\begin{theorem}}{\end{theorem}}
\newenvironment{yinli}{\begin{lemma}}{\end{lemma}}
\newenvironment{mingti}{\begin{proposition}}{\end{proposition}}
\newenvironment{tuilun}{\begin{corollary}}{\end{corollary}}
\newenvironment{dingyi}{\begin{definition}}{\end{definition}}
\newenvironment{li}{\begin{example}}{\end{example}}
\newenvironment{zhu}{\begin{remark}}{\end{remark}}
\newenvironment{biji}{\begin{note}}{\end{note}}
\newenvironment{kaoshi}{\begin{examnote}}{\end{examnote}}

\renewcommand{\qedsymbol}{\ensuremath{\square}}

% PDF sidebar outline + printed TOC for QE exam transcripts.
% Load in the preamble AFTER hyperref and AFTER \newcommand{\examstatement}.
\hypersetup{
  unicode=true,
  bookmarks=true,
  bookmarksopen=true,
  bookmarksopenlevel=3,
  bookmarksnumbered=false,
  bookmarksdepth=3
}
\IfFileExists{bookmark.sty}{%
  \usepackage{bookmark}%
  \bookmarksetup{open,openlevel=3,numbered=false,depth=3}%
}{}
% Unnumbered in print, but still written to the TOC / PDF outline
% down to \subsubsection. Starred headings create neither.
\setcounter{secnumdepth}{-1}
\setcounter{tocdepth}{3}
\makeatletter
\@ifundefined{examstatement}{}{%
  \let\qe@origexamstatement\examstatement
  \renewcommand{\examstatement}[1]{%
    \par\phantomsection
    \addcontentsline{toc}{subsection}{#1}%
    \qe@origexamstatement{#1}%
  }%
}
\makeatother


\title{\textbf{QUALIFYING EXAM: PROBABILITY \& STATISTICS}\\[0.5em]\Large FALL 2023}
\date{}
\author{}

\begin{document}

\maketitle
\vspace{-3em}

\tableofcontents

\section{Instructions}
\begin{itemize}
    \item 本试卷共3页, 8道大题, 总分为100分.
    \item 考生默认遵守考试纪律, 不遵守者后果自负.
    \item 所有的解答请写出必要的细节,推理依据和推理过程.
\end{itemize}

\section{Statement of the problems}

\examstatement{Problem 1. [10 pts.]}
Let \( (X, Y) \) be a random vector of joint density

\[
f(x, y) = \frac{1}{2\pi \sqrt{1 - \rho^2}} \exp\left( -\frac{x^2 - 2\rho xy + y^2}{2(1 - \rho^2)} \right).
\]

\begin{enumerate}[label=(\alph*)]
    \item Prove that \( f \) is a probability density. Calculate the marginal distributions of \( X \) and \( Y \). Determine the condition of \( \rho \) such that \( X \) and \( Y \) are independent.
    \item Define \( R = \sqrt{X^2 + Y^2} \) and \( \Phi \in [0, 2\pi] \) by

    \[
    \cos \Phi = \frac{X}{R}, \qquad \sin \Phi = \frac{Y}{R}.
    \]

    Determine the law of \( (R, \Phi) \) and then the law of \( \Phi \).
\end{enumerate}

\examstatement{Problem 2. [10 pts.]}
Let \( (X_n)_{n \in \mathbb{N}_+} \) be i.i.d.\ random variables of density

\[
f(x) = \frac{1}{2} e^{-|x|}, \qquad x \in \mathbb{R},
\]

and we denote by \( \mathcal{F}_n \) its natural filtration \( \mathcal{F}_n = \sigma((X_i)_{1 \le i \le n}) \). We define \( M_n(\theta) \) as

\[
M_n(\theta) := \exp\left( \theta \sum_{k=1}^n X_k - A_n(\theta) \right).
\]

Determine \( (A_n(\theta))_{n \in \mathbb{N}_+} \) such that for any \( |\theta| < 1 \), the sequence \( (M_n(\theta))_{n \in \mathbb{N}_+} \) is a martingale with respect to the natural filtration \( (\mathcal{F}_n)_{n \in \mathbb{N}_+} \).

\examstatement{Problem 3. [10 pts.]}
There are \( N \) identical coins and every coin has two sides, which are denoted by ``Head'' and ``Tail''. Every round we draw a coin uniformly and flip it. Suppose that every draw is independent and let \( X_n \) be the number of heads after \( n \)-th round.
\begin{enumerate}[label=(\alph*)]
    \item Prove that \( (X_n)_{n \in \mathbb{N}_+} \) is a Markov chain and write down the transition probabilities.
    \item Prove that there exists a unique stationary distribution \( \mu \) for this Markov chain. Determine the law of \( \mu \).
\end{enumerate}

\examstatement{Problem 4. [10 pts.]}
Alice has 3 identical dices. Every dice has faces marked \( \{1, 2, 3, 4, 5, 6\} \) and can sample value uniformly from these faces. Alice plays the following game: in every round, she can choose any subset of dices and resample their values, while keep the other dices as last round. All the samplings of dices are independent, and Alice's decision does not change the law of dices. Let \( T \) be the number of rounds that Alice realizes the 3 dices of configuration ``6, 6, 6'' for the first time.
\begin{enumerate}[label=(\alph*)]
    \item What is the optimal strategy of Alice to minimize the waiting time \( T \)?
    \item Under the optimal strategy, calculate the expectation of \( T \).
\end{enumerate}

\examstatement{Problem 5. [15 pts.]}
Consider the Bernoulli percolation on an infinite \( d \)-regular (\( d \in \mathbb{N} \) and \( d \ge 3 \)) tree \( T \) with root \( o \). More precisely, \( T \) is an infinite tree where each vertex has \( d \) neighborhoods. Each edge is open with probability \( p \), and closed with probability \( 1 - p \), independently of the states of other edges. Let \( \mathcal{C}(o) \) be the connected component (called \emph{cluster}) containing \( o \), i.e., \( \mathcal{C}(o) \) is the subset of all vertices that can be reached from \( o \) through open edges. Let \( \theta(p) \) be the probability that \( o \) is in an infinite cluster, i.e., \( \theta(p) = \mathbb{P}(|\mathcal{C}(o)| = \infty) \).
\begin{enumerate}[label=(\alph*)]
    \item Prove that there exists a \( p_c \in (0, 1) \) such that if \( p > p_c \), then \( \theta(p_c) > 0 \), while if \( p < p_c \), then \( \theta(p_c) = 0 \). Give the value of \( p_c \).
    \item Fix any \( p \in (0, p_c) \), prove that there exist constants \( c, C > 0 \) (which may depend on \( p_c \) but do not depend on \( k \)) such that for all \( k \ge 2 \), there is

    \[
    \mathbb{P}(|\mathcal{C}(o)| \ge k) \le C \exp(-ck).
    \]
\end{enumerate}

\begin{center}
\small
图~1: An illustration of a 3-regular tree.
\end{center}

\examstatement{Problem 6. [15 pts.]}
Let \( X \) be an \( n \times n \) random matrix with i.i.d.\ standard Gaussian entries, i.e., the entries \( X_{ij} \) are independent standard normal random variables.
\begin{enumerate}[label=(\alph*)]
    \item Given any deterministic orthogonal matrices \( U \) and \( V \), prove that \( UXV \) has the same law as \( X \).
    \item Given any deterministic unit vector \( u \in \mathbb{R}^n \), calculate the mean and variance of \( Y_n = u^T X^T X u \), and prove the central limit theorem for \( Y_n \) as \( n \to \infty \) (i.e., \( Y_n \), after proper shift and rescaling, converges in law to the standard normal distribution).
    \item Given any deterministic unit vectors \( u, v \in \mathbb{R}^n \), calculate the mean and variance of \( Z_n = u^T X^T X u + \sqrt{n}\, v^T X v \), and prove the central limit theorem for \( Z_n \) as \( n \to \infty \). (Hint: You may try the moment method or Stein's method, i.e., a random variable \( x_n \) converges in law to the standard normal distribution if for any differentiable function \( f \) with bounded derivatives, \( \mathbb{E}(f'(x_n) - x_n f(x_n)) \to 0 \) as \( n \to \infty \).)
\end{enumerate}

\examstatement{Problem 7. [15 pts.]}
Assume that \( X_1, \dots, X_n \) are i.i.d.\ \( N(\mu, \sigma^2) \) and consider testing the hypotheses \( H_0 : \mu = \mu_0 \), \( \sigma^2 = \sigma_0^2 \) versus \( H_1 : \mu = \mu_1 \), \( \sigma^2 = \sigma_1^2 \). Note that the density for \( N(\mu, \sigma^2) \) is

\[
\frac{1}{\sqrt{2\pi}\,\sigma} e^{-\frac{(x-\mu)^2}{2\sigma^2}}, \qquad -\infty < x < \infty,\quad -\infty < \mu < \infty,\quad \sigma > 0.
\]

\begin{enumerate}[label=(\alph*)]
    \item Find the uniformly most powerful (UMP) size \( \alpha \) test of testing \( H_0 \) vs \( H_1 \) if \( \mu_0 < \mu_1 \) and \( \sigma_0^2 = \sigma_1^2 \).
    \item Find the UMP size \( \alpha \) test of testing \( H_0 \) vs \( H_1 \) if \( \mu_0 = \mu_1 \) and \( \sigma_0^2 < \sigma_1^2 \).
    \item Assume \( \sigma^2 = 1 \), find the UMP level \( \alpha \) test of testing \( H_{0}^{0} : \mu \ge 1 \) versus \( H_{1}^{0} : \mu < 1 \).
\end{enumerate}

\examstatement{Problem 8. [15 pts.]}
Consider the following distribution (Cauchy) \( f(x \mid \theta) = \frac{1}{\pi[1+(x-\theta)^2]} \), \( -\infty < x < \infty \), \( -\infty < \theta < \infty \).
\begin{enumerate}[label=(\alph*)]
    \item Let \( X_1, \dots, X_n \) be a random sample from the above distribution, find a minimal sufficient statistic for \( \theta \).
    \item Let \( X_1 \) and \( X_2 \) be i.i.d with the above density. Let \( x_1 \) and \( x_2 \) be the observations and set \( \Delta = \frac{1}{2}(x_1 - x_2) \). Let \( \hat{\theta} \) denote the maximum likelihood estimator (MLE). Show that if \( |\Delta| \le 1 \), then the MLE exists and is unique. Give the MLE when \( |\Delta| \le 1 \).
    \item Show that if \( |\Delta| > 1 \), then the MLE is not unique. Find the values of \( \theta \) that maximize the likelihood when \( |\Delta| > 1 \).
\end{enumerate}

\noindent\rule{\linewidth}{0.4pt}

\section{Statement and solutions}

\subsection{Problem 1. [10 pts.]}
Let \( (X, Y) \) be a random vector of joint density

\[
f(x, y) = \frac{1}{2\pi \sqrt{1 - \rho^2}} \exp\left( -\frac{x^2 - 2\rho xy + y^2}{2(1 - \rho^2)} \right).
\]

\begin{enumerate}[label=(\alph*)]
    \item Prove that \( f \) is a probability density. Calculate the marginal distributions of \( X \) and \( Y \). Determine the condition of \( \rho \) such that \( X \) and \( Y \) are independent.
    \item Define \( R = \sqrt{X^2 + Y^2} \) and \( \Phi \in [0, 2\pi] \) by

    \[
    \cos \Phi = \frac{X}{R}, \qquad \sin \Phi = \frac{Y}{R}.
    \]

    Determine the law of \( (R, \Phi) \) and then the law of \( \Phi \).
\end{enumerate}

\subsection{Problem 2. [10 pts.]}
Let \( (X_n)_{n \in \mathbb{N}_+} \) be i.i.d.\ random variables of density

\[
f(x) = \frac{1}{2} e^{-|x|}, \qquad x \in \mathbb{R},
\]

and we denote by \( \mathcal{F}_n \) its natural filtration \( \mathcal{F}_n = \sigma((X_i)_{1 \le i \le n}) \). We define \( M_n(\theta) \) as

\[
M_n(\theta) := \exp\left( \theta \sum_{k=1}^n X_k - A_n(\theta) \right).
\]

Determine \( (A_n(\theta))_{n \in \mathbb{N}_+} \) such that for any \( |\theta| < 1 \), the sequence \( (M_n(\theta))_{n \in \mathbb{N}_+} \) is a martingale with respect to the natural filtration \( (\mathcal{F}_n)_{n \in \mathbb{N}_+} \).

\subsection{Problem 3. [10 pts.]}
There are \( N \) identical coins and every coin has two sides, which are denoted by ``Head'' and ``Tail''. Every round we draw a coin uniformly and flip it. Suppose that every draw is independent and let \( X_n \) be the number of heads after \( n \)-th round.
\begin{enumerate}[label=(\alph*)]
    \item Prove that \( (X_n)_{n \in \mathbb{N}_+} \) is a Markov chain and write down the transition probabilities.
    \item Prove that there exists a unique stationary distribution \( \mu \) for this Markov chain. Determine the law of \( \mu \).
\end{enumerate}

\subsection{Problem 4. [10 pts.]}
Alice has 3 identical dices. Every dice has faces marked \( \{1, 2, 3, 4, 5, 6\} \) and can sample value uniformly from these faces. Alice plays the following game: in every round, she can choose any subset of dices and resample their values, while keep the other dices as last round. All the samplings of dices are independent, and Alice's decision does not change the law of dices. Let \( T \) be the number of rounds that Alice realizes the 3 dices of configuration ``6, 6, 6'' for the first time.
\begin{enumerate}[label=(\alph*)]
    \item What is the optimal strategy of Alice to minimize the waiting time \( T \)?
    \item Under the optimal strategy, calculate the expectation of \( T \).
\end{enumerate}

\subsection{Problem 5. [15 pts.]}
Consider the Bernoulli percolation on an infinite \( d \)-regular (\( d \in \mathbb{N} \) and \( d \ge 3 \)) tree \( T \) with root \( o \). More precisely, \( T \) is an infinite tree where each vertex has \( d \) neighborhoods. Each edge is open with probability \( p \), and closed with probability \( 1 - p \), independently of the states of other edges. Let \( \mathcal{C}(o) \) be the connected component (called \emph{cluster}) containing \( o \), i.e., \( \mathcal{C}(o) \) is the subset of all vertices that can be reached from \( o \) through open edges. Let \( \theta(p) \) be the probability that \( o \) is in an infinite cluster, i.e., \( \theta(p) = \mathbb{P}(|\mathcal{C}(o)| = \infty) \).
\begin{enumerate}[label=(\alph*)]
    \item Prove that there exists a \( p_c \in (0, 1) \) such that if \( p > p_c \), then \( \theta(p_c) > 0 \), while if \( p < p_c \), then \( \theta(p_c) = 0 \). Give the value of \( p_c \).
    \item Fix any \( p \in (0, p_c) \), prove that there exist constants \( c, C > 0 \) (which may depend on \( p_c \) but do not depend on \( k \)) such that for all \( k \ge 2 \), there is

    \[
    \mathbb{P}(|\mathcal{C}(o)| \ge k) \le C \exp(-ck).
    \]
\end{enumerate}

\begin{center}
\small
图~1: An illustration of a 3-regular tree.
\end{center}

\begin{remark*}
The official wording in (a) writes \( \theta(p_c) \); the intended dichotomy is \( \theta(p) > 0 \) for \( p > p_c \) and \( \theta(p) = 0 \) for \( p < p_c \).
\end{remark*}

\subsection{Problem 6. [15 pts.]}
Let \( X \) be an \( n \times n \) random matrix with i.i.d.\ standard Gaussian entries, i.e., the entries \( X_{ij} \) are independent standard normal random variables.
\begin{enumerate}[label=(\alph*)]
    \item Given any deterministic orthogonal matrices \( U \) and \( V \), prove that \( UXV \) has the same law as \( X \).
    \item Given any deterministic unit vector \( u \in \mathbb{R}^n \), calculate the mean and variance of \( Y_n = u^T X^T X u \), and prove the central limit theorem for \( Y_n \) as \( n \to \infty \) (i.e., \( Y_n \), after proper shift and rescaling, converges in law to the standard normal distribution).
    \item Given any deterministic unit vectors \( u, v \in \mathbb{R}^n \), calculate the mean and variance of \( Z_n = u^T X^T X u + \sqrt{n}\, v^T X v \), and prove the central limit theorem for \( Z_n \) as \( n \to \infty \). (Hint: You may try the moment method or Stein's method, i.e., a random variable \( x_n \) converges in law to the standard normal distribution if for any differentiable function \( f \) with bounded derivatives, \( \mathbb{E}(f'(x_n) - x_n f(x_n)) \to 0 \) as \( n \to \infty \).)
\end{enumerate}

\subsection{Problem 7. [15 pts.]}
Assume that \( X_1, \dots, X_n \) are i.i.d.\ \( N(\mu, \sigma^2) \) and consider testing the hypotheses \( H_0 : \mu = \mu_0 \), \( \sigma^2 = \sigma_0^2 \) versus \( H_1 : \mu = \mu_1 \), \( \sigma^2 = \sigma_1^2 \). Note that the density for \( N(\mu, \sigma^2) \) is

\[
\frac{1}{\sqrt{2\pi}\,\sigma} e^{-\frac{(x-\mu)^2}{2\sigma^2}}, \qquad -\infty < x < \infty,\quad -\infty < \mu < \infty,\quad \sigma > 0.
\]

\begin{enumerate}[label=(\alph*)]
    \item Find the uniformly most powerful (UMP) size \( \alpha \) test of testing \( H_0 \) vs \( H_1 \) if \( \mu_0 < \mu_1 \) and \( \sigma_0^2 = \sigma_1^2 \).
    \item Find the UMP size \( \alpha \) test of testing \( H_0 \) vs \( H_1 \) if \( \mu_0 = \mu_1 \) and \( \sigma_0^2 < \sigma_1^2 \).
    \item Assume \( \sigma^2 = 1 \), find the UMP level \( \alpha \) test of testing \( H_{0}^{0} : \mu \ge 1 \) versus \( H_{1}^{0} : \mu < 1 \).
\end{enumerate}

\subsection{Problem 8. [15 pts.]}
Consider the following distribution (Cauchy) \( f(x \mid \theta) = \frac{1}{\pi[1+(x-\theta)^2]} \), \( -\infty < x < \infty \), \( -\infty < \theta < \infty \).
\begin{enumerate}[label=(\alph*)]
    \item Let \( X_1, \dots, X_n \) be a random sample from the above distribution, find a minimal sufficient statistic for \( \theta \).
    \item Let \( X_1 \) and \( X_2 \) be i.i.d with the above density. Let \( x_1 \) and \( x_2 \) be the observations and set \( \Delta = \frac{1}{2}(x_1 - x_2) \). Let \( \hat{\theta} \) denote the maximum likelihood estimator (MLE). Show that if \( |\Delta| \le 1 \), then the MLE exists and is unique. Give the MLE when \( |\Delta| \le 1 \).
    \item Show that if \( |\Delta| > 1 \), then the MLE is not unique. Find the values of \( \theta \) that maximize the likelihood when \( |\Delta| > 1 \).
\end{enumerate}
\end{document}
